(1) 镜面向外位移 x，光在该臂往返传播，多走的光路是 \(2x\)，所以光程差变化 \(\Delta=2x\)。对应相位变化
\[
\delta\phi=\frac{2\pi}{\lambda}\Delta=\frac{4\pi x}{\lambda}.
\]

(2) 若 x(t)=vt，探测器两束光的相位差为 \(\phi(t)=\phi(0)+4\pi vt/\lambda\)，故条纹相位变化率 \(d\phi/dt=4\pi v/\lambda\)。每增加 \(2\pi\) 相位恰出现一个干涉条纹周期。

(3) 条纹通过率是
\[
f_{\rm beat}=\frac1{2\pi}\left|\frac{d\phi}{dt}\right|
=\frac{2|v|}{\lambda}.
\]
单个反射光在缓慢移动镜面上的多普勒频移一阶为 \(\Delta\nu\simeq-2\nu v/c=-2v/\lambda\)（镜面远离时降频）。与固定臂的参考波相干产生拍频 \(|\Delta\nu|=2|v|/\lambda\)，独立验证。

(4) 经过 N 个完整条纹对应 \(|\delta\phi|=2\pi N\)，所以
\[
|\Delta x|=N\lambda/2.
\]
当 v=0 时相位恒定，条纹不流动，拍频为0。本结论属于低速一级近似；接近光速时需用精确相对论双程多普勒公式。
